% ---
% titre: Révisions d'examen
% theme: NO
% chapitre: Nombres rationnels
% sous_chapitre: Résoudre des problèmes
% niveau: [11e]
% type: EVAL
% tags: [c-soustraction,c-multiplication,c-fractions,c-operations,c-probleme,p-question-TS]
% difficulte: 1
% corrige: true
% ---

\question[5]
Un élève a 90 heures de travail à disposition pour réviser ses trois branches d'examen. Il utilise $\dfrac{1}{5}$ du temps total pour la première branche, $\dfrac{7}{10}$ pour la deuxième branche et le reste pour la troisième branche.

\begin{parts}
	\vspace{20pt}
	\part
		Quelle fraction du temps total utilise-il pour la troisième branche?
		
	\vspace{10pt}{\footnotesize\raggedleft Espace pour ta démarche et tes calculs\par}
	\begin{solutionorgrid}[120pt]
	
		\vspace{10pt}
		$\dfrac{10}{10} - \dfrac{1}{5} - \dfrac{7}{10} \textbf{\,(0.5pt)} = \dfrac{10}{10} - \dfrac{2}{10} - \dfrac{7}{10} \textbf{\,(0.5pt)} = \dfrac{1}{10} \textbf{\,(0.5pt)}$
	\end{solutionorgrid}
	
	\begin{tcolorbox}[enhanced jigsaw, interior hidden, colframe=box, parbox=false]%
	Réponse: 
	\sol{\vspace{20pt}}{Il utilise \textbf{$\dfrac{1}{10}$} du temps pour réviser la troisième branche. \textbf{\,(0.5pt)}}
	\end{tcolorbox}	
		
	\vspace{20pt}
	\part
		Calcule le nombre d'heures passées à réviser pour chacune des branches. 
		
	\vspace{10pt}{\footnotesize\raggedleft Espace pour ta démarche et tes calculs\par}
	\begin{solutionorgrid}[200pt]
	
		\vspace{10pt}
		Branche 1: 
		$\dfrac{1}{5} \cdot 90 \textbf{\,(1pt)} = 18\,h \textbf{\,(0.5pt)}$
		
		\vspace{10pt}
		Branche 2: 
		$\dfrac{7}{10} \cdot 90 = 63\,h \textbf{\,(0.5pt)}$
		
		\vspace{10pt}
		Branche 3: 
		$\dfrac{1}{10} \cdot 90 = 9\,h \textbf{\,(0.5pt)}$
	\end{solutionorgrid}
			
	\begin{tcolorbox}[enhanced jigsaw, interior hidden, colframe=box, parbox=false]%
	Réponse: 
	\sol{\vspace{20pt}}{Il passe \textbf{18\,h} à réviser la première branche, \textbf{63\,h} à réviser la deuxième branche et \textbf{9\,h} à réviser la troisième branche. \textbf{\,(0.5pt)}}
	\end{tcolorbox}
\end{parts}


